\documentclass[11pt]{article}

\usepackage[utf8]{inputenc}
\usepackage[T1]{fontenc}
\usepackage[margin=1in]{geometry}
\usepackage{amsmath,amssymb,mathtools}
\usepackage{booktabs}
\usepackage{array}
\usepackage{enumitem}
\usepackage[hidelinks]{hyperref}

\DeclareMathOperator{\score}{score}
\DeclareMathOperator{\Bin}{Bin}
\DeclarePairedDelimiter{\abs}{\lvert}{\rvert}
\DeclarePairedDelimiter{\ceil}{\lceil}{\rceil}
\newcommand{\W}{\mathcal{W}}
\newcommand{\A}{\mathcal{A}}
\newcommand{\klo}{k_{\mathrm{lo}}}
\newcommand{\khi}{k_{\mathrm{hi}}}
\newcommand{\kaplo}{\kappa_{\mathrm{lo}}}
\newcommand{\kaphi}{\kappa_{\mathrm{hi}}}
\newcommand{\Sseg}{S_{\mathrm{seg}}}

\newcounter{prob}
\newcommand{\problem}[1]{\refstepcounter{prob}\medskip\noindent\textbf{Problem \theprob} (#1)\textbf{.}\ }
\newcommand{\easy}{$\star$}
\newcommand{\medium}{$\star\star$}
\newcommand{\hard}{$\star\star\star$}

\title{Proof Problems for the Alignment Kernel}
\author{TB Hardware Accelerator project}
\date{}

\begin{document}
\maketitle

\noindent
Our project lines up two strings of DNA letters to find where they differ. The
standard method fills in a whole table. Our design first guesses where the
answer lies (the \emph{vote}) and then fills in only a narrow strip of the table
(the \emph{window}). The statements below say how much work this saves and when
it is guaranteed to give the same answer.

We have draft proofs, and nobody has checked them yet. \textbf{Please prove
these by hand, without looking at the drafts.} Afterwards we will compare. If
you find a statement that is false or needs an extra condition, that is the
most useful result of all.

\textbf{Take a whole part, not single problems.} The problems inside a part
build on each other: the early ones are the steps and the last one is the goal.
Section~\ref{sec:map} shows how the parts connect.

Difficulty: \easy\ counting or a few lines, \medium\ a real argument, \hard\
longer, usually induction. Write up each part as one piece, with the part
letter and your names on it.

\section{The rules}

\begin{itemize}[nosep]
\item The \emph{read} is $a = a_1 \dots a_m$ and the \emph{reference} is
      $b = b_1 \dots b_n$. Letters are A, C, G, T.
\item Scores: a match scores $M > 0$, a mismatch scores $x \le M$, and each
      gapped letter scores $g < 0$. The project uses $M = 2$, $x = -1$,
      $g = -2$, $m = 150$, $n \le 180$.
\item $s(a_i,b_j) = M$ if $a_i = b_j$, and $x$ otherwise.
\item \emph{Nodes} are pairs $(i,j)$ with $0 \le i \le m$ and $0 \le j \le n$.
      A node with $i,j \ge 1$ is a \emph{cell}. The \emph{diagonal} of a node
      is $j - i$.
\item A \emph{path} is a sequence of steps. A diagonal step
      $(i-1,j-1)\to(i,j)$ scores $s(a_i,b_j)$. A down step $(i-1,j)\to(i,j)$
      and a right step $(i,j-1)\to(i,j)$ each score $g$; these are \emph{gap
      steps}. The score of a path is the sum of its steps. The empty path
      scores $0$.
\end{itemize}

\paragraph{The table.} $H(i,0) = H(0,j) = 0$, and for $i,j \ge 1$
\begin{equation}\label{eq:sw}
H(i,j) = \max\bigl\{\,0,\; H(i-1,j-1) + s(a_i,b_j),\; H(i-1,j) + g,\; H(i,j-1) + g \,\bigr\}.
\end{equation}
The answer is $S = \max_{i,j} H(i,j)$.

\paragraph{A small example.} Reference GATTACA, read GATACA, with $M=2$,
$x=-1$, $g=-2$:
\begin{center}
\begin{tabular}{c|rrrrrrr}
  & G & A & T & T & A & C & A \\ \hline
G & 2 & 0 & 0 & 0 & 0 & 0 & 0 \\
A & 0 & 4 & 2 & 0 & 2 & 0 & 2 \\
T & 0 & 2 & 6 & 4 & 2 & 1 & 0 \\
A & 0 & 2 & 4 & 5 & 6 & 4 & 3 \\
C & 0 & 0 & 2 & 3 & 4 & 8 & 6 \\
A & 0 & 2 & 0 & 1 & 5 & 6 & 10
\end{tabular}
\end{center}
Here $S = 10$. The best path uses one gap step: the read is missing one T.

\section{How the parts connect}\label{sec:map}

\begin{center}
\renewcommand{\arraystretch}{1.25}
\begin{tabular}{@{}cp{0.33\textwidth}p{0.2\textwidth}p{0.22\textwidth}c@{}}
\toprule
Part & What it establishes & Problems & Needs & Level \\
\midrule
A & What the numbers in the table mean, and how big they can get & 1--4 & only the rules & \hard \\
B & The score rewritten as a cost & 5 & only the rules & \medium \\
C & When the window gives the same answer as the full table & 6--10 & Problem 1 from Part A & \medium \\
D & When the vote is guaranteed to pick a good place & 11--15 & Problems 7 and 8 from Part C & \hard \\
E & How much hardware and time each design needs & 16 & nothing & \easy \\
F & Two separate bounds & 17--18 & nothing & \medium \\
\bottomrule
\end{tabular}
\end{center}

\noindent
Parts A, B, E and F can each be done alone. Part C should be done by people who
have read Problem 1. Part D should be done after Part C, or by a group that is
willing to take Problems 7 and 8 as given. If you use a result from another
part, quote it by its problem number; you do not have to prove it again.

\section{Part A: the full table}

\noindent\textbf{Level: \hard.} The hardest step is Problem 3, an induction
with four cases. Skills: induction, careful case analysis. Rough time: 2 to 3
hours.

\noindent\textbf{Goal.} Problems 2 and 4: how large a score can be, and how few
bits a cell needs. \textbf{Chain.} $1 \to 2$ and $3 \to 4$. Problem 1 says what
$H$ means; Problem 3 is the long induction.

\problem{\medium}
Prove that $H(i,j)$ is the maximum score over all paths that end at the node
$(i,j)$, the empty path included.

\problem{\easy}
Prove that $0 \le H(i,j) \le M\min(i,j)$. Deduce that $S \le M\min(m,n)$ and
that $S$ fits in $\ceil{\log_2(M\min(m,n)+1)}$ bits. Evaluate this for the
project values.

\problem{\hard}
Prove that for all $1 \le i \le m$ and $1 \le j \le n$:
\begin{align*}
g &\le H(i,j) - H(i,j-1) \le M - g, \\
g &\le H(i,j) - H(i-1,j) \le M - g, \\
x &\le H(i,j) - H(i-1,j-1) \le M.
\end{align*}

\problem{\easy}
Using Problem 3, how many different values can a side-by-side difference take?
A same-diagonal difference? How many bits does each need? Evaluate for the
project values.

\section{Part B: score as a cost}

For a path $P$ from $(i_0,j_0)$ to $(i_1,j_1)$ let $\mu$, $\nu$, $\gamma$ be its
numbers of matches, mismatches and gap steps, and let
$u = i_0 + (m-i_1) + j_0 + (n-j_1)$ be the number of letters it leaves
unaligned. Define
\[
c(P) = (M-x)\,\nu + \Bigl(\frac{M}{2}-g\Bigr)\gamma + \frac{M}{2}\,u .
\]

\noindent\textbf{Level: \medium.} One identity. Skills: algebra and keeping
track of what each step of a path uses up. Rough time: 30 minutes.

\noindent\textbf{Goal.} Problem 5, a single identity. It uses nothing except
the rules.

\problem{\medium}
Prove that $\score(P) = \frac{M}{2}(m+n) - c(P)$ for every path $P$. Deduce
that $S = \frac{M}{2}(m+n) - \min_P c(P)$. What are the three cost
coefficients for the project values?

\section{Part C: the window}

Fix an integer $d_0$ and a half-width $w \ge 0$, and let $W = 2w+1$. The
\emph{window} is the set of cells
\[
\W = \bigl\{(i,j) : 1\le i\le m,\ 1\le j\le n,\ \abs{j-i-d_0}\le w\bigr\}.
\]
$H_{\W}$ is defined by the recurrence \eqref{eq:sw} on the cells of $\W$ and is
$0$ at every other node. Let $S_{\W} = \max H_{\W}$. A path is
\emph{admissible} if every node on it, except possibly its start node, is a
cell of $\W$.

\noindent\textbf{Level: \medium.} The hardest step is Problem 7, an induction.
The rest are short once 7 is done. Skills: induction, reasoning about subsets.
Rough time: 1.5 to 2 hours.

\noindent\textbf{Goal.} Problem 10: if the best path fits in the window, the
window gives the same answer. \textbf{Chain.} Problem 6 stands alone and says
how much work is saved. Then $7 \to 8$, and $7, 8, 9 \to 10$. Problem 7 is the
window version of Problem 1 and can follow the same pattern.

\problem{\easy}
Prove that $\abs{\W} \le mW$. What fraction of the full table is this for
$m=150$, $n=180$, $w=8$?

\problem{\medium}
Prove that for a cell $(i,j)\in\W$, $H_{\W}(i,j)$ is the maximum score over
admissible paths ending at $(i,j)$, the empty path included.

\problem{\easy}
Prove that $S_{\W} \le S$.

\problem{\medium}
Prove that the diagonals of the nodes of a path form a set of consecutive
integers $\{\klo,\dots,\khi\}$, and that $\khi - \klo$ is at most the number of
gap steps of the path.

\problem{\medium}
Suppose $S > 0$ and some path $P^{*}$ with $\score(P^{*}) = S$ has diagonal
range $[\klo,\khi]$ with $\khi - w \le d_0 \le \klo + w$. Prove that
$S_{\W} = S$. Show that the condition holds whenever $d_0$ is one of the
diagonals of $P^{*}$ and $P^{*}$ has at most $w$ gap steps.

\section{Part D: the vote}

In this part indices start at $0$: the read is $r_0\dots r_{m-1}$, the window
of the reference is $b_0\dots b_{n-1}$, and diagonal $k$ pairs $r_i$ with
$b_{i+k}$. Fix $\lambda \ge 0$. For $-(m-1)\le k\le n-1$ let
\begin{align*}
I(k) &= \{\,i : 0\le i<m,\ 0\le i+k<n\,\}, & N(k) &= \abs{I(k)},\\
V_r(k) &= \abs[\big]{\{\,i\in I(k) : r_i = b_{i+k}\,\}}, &
Z_r(k) &= (1+\lambda)V_r(k) - \lambda N(k).
\end{align*}
The vote returns a diagonal $d_0$ that maximises $Z_r$.

\paragraph{Read model.} Let $E$ be the genome around the window, with
$b_j = E[\omega+j]$. A \emph{template} $t_0\dots t_{m-1}$ has a start $p$ and
at most one indel of length $g$ at read position $q$:
\begin{itemize}[nosep]
\item no indel: $t_i = E[p+i]$;
\item deletion: $t_i = E[p+i]$ for $i<q$, and $t_i = E[p+g+i]$ for $i\ge q$;
\item insertion: $t_i = E[p+i]$ for $i<q$; $t_i$ is unconstrained for
      $q\le i<q+g$; and $t_i = E[p-g+i]$ for $i\ge q+g$.
\end{itemize}
$F$ is the set of unconstrained positions. A read $r$ \emph{follows the
template with at most $e$ substitutions} if $r_i = t_i$ for all $i\notin F$
except at most $e$ positions.

The \emph{true diagonals} are $k_1 = p-\omega$ and $k_2 = k_1$ (no indel),
$k_1+g$ (deletion) or $k_1-g$ (insertion). Put $\kaplo = \min(k_1,k_2)$,
$\kaphi = \max(k_1,k_2)$. The \emph{true alignment} $T$ pairs each read letter
with the reference letter it was copied from. The set of \emph{acceptable}
diagonals is $\A = \{d : \kaphi - w \le d \le \kaplo + w\}$. Define
\[
V^{-}(k) = \abs[\big]{\{\,i\in I(k)\setminus F : t_i = b_{i+k}\,\}},\qquad
V^{+}(k) = V^{-}(k) + \abs{I(k)\cap F},
\]
$Z^{\pm}(k) = (1+\lambda)V^{\pm}(k) - \lambda N(k)$, and the \emph{margin}
\[
\Delta_\lambda = \max_{k\in\A} Z^{-}(k) - \max_{k\notin\A} Z^{+}(k).
\]
Here $k$ ranges over $-(m-1),\dots,n-1$. We assume $g \le 2w$, so that $\A$ is
not empty. If every diagonal is acceptable, the second maximum is over an empty
set and is taken to be $-\infty$.

\noindent\textbf{Level: \hard.} The proofs are short; the difficulty is the
amount of notation, and Problem 12 needs a chain of four inequalities in the
right order. Skills: reading definitions precisely, inequalities, a little
probability for Problem 14. Rough time: 2 to 3 hours.

\noindent\textbf{Goal.} Problem 13: if the read has few enough substitutions,
the window chosen by the vote contains the true alignment. \textbf{Chain.}
$11 \to 12 \to 13$, and Problem 13 also uses Problems 7 and 8. Problem 14 follows
from 12. Problem 15 is a separate count. Read the definitions above twice
before starting; they are the hard part.

\problem{\medium}
Prove that if $r$ follows the template with at most $e$ substitutions, then
$V^{-}(k) - e \le V_r(k) \le V^{+}(k) + e$ for every $k$.

\problem{\hard}
Prove: if $r$ follows the template with at most $e$ substitutions and
$2(1+\lambda)e < \Delta_\lambda$, then every maximiser of $Z_r$ is acceptable.

\problem{\medium}
Under the hypothesis of Problem 12, prove that the window centred on $d_0$
contains every diagonal of the true alignment $T$, and that
$\score(T) \le S_{\W} \le S$.

\problem{\medium}
Suppose each read letter is substituted independently with probability
$\varepsilon$, and the template tolerates $e^{*}$ substitutions in the sense of
Problem 12. Prove that the probability that the vote is not acceptable is at
most $\Pr[\Bin(m,\varepsilon) > e^{*}]$. Evaluate it for $m=150$, $e^{*}=19$,
$\varepsilon = 0.02$.

\problem{\easy}
Prove that computing $V_r(k)$ for every $k$ by direct comparison uses exactly
$mn$ comparisons of letters.

\section{Part E: counting the hardware}

A machine that fills the full table one anti-diagonal per clock tick uses
$\min(m,n)$ cells and $m+n-1$ ticks, and stores $2mn$ bits of direction
pointers. The window machine uses $W$ cells and $m$ ticks and stores $2mW$
bits. The vote takes $m+n-1$ ticks.

\noindent\textbf{Level: \easy.} Arithmetic and one inequality. A good place to
start. Rough time: 20 minutes.

\noindent\textbf{Goal.} Problem 16, which is counting only. Part (c) is the
point of it.

\problem{\easy}
(a) Evaluate all of these for $m=150$, $n=180$, $W=17$.
(b) A read goes through the vote and then the window. Is that more or fewer
ticks than the full-table machine? Prove your answer for all $m,n\ge1$.
(c) State precisely what the new design saves and what it does not.

\section{Part F: two side results}

\noindent\textbf{Level: \medium.} Problem 17 needs one neat inequality about
three numbers; Problem 18 is one line. Rough time: 45 minutes.

\noindent Problems 17 and 18 do not depend on each other or on anything else.

\problem{\medium}
Split the read into consecutive segments of length at most $\Sseg$. Segment $t$
votes for a diagonal $d_t$. A row in segment $t$ computes the diagonals from
$\min(d_{t-1},d_t,d_{t+1})-w$ to $\max(d_{t-1},d_t,d_{t+1})+w$ (a missing
neighbour is left out). Let $G = \sum_t\abs{d_{t+1}-d_t}$. Prove that at most
$mW + 2\Sseg G$ cells are computed.

\problem{\easy}
A read skips the window with probability $\rho$. The vote costs
$C_{\mathrm{vote}}$ and the window costs $C_{\mathrm{SW}}$. Prove that the
expected cost per read is $C_{\mathrm{vote}} + (1-\rho)\,C_{\mathrm{SW}}$.

\section{Part G: open questions}

\noindent\textbf{Level: unknown.} We do not have proofs of these. We do not
know whether they are true. Q1 and Q4 look the hardest.

\begin{enumerate}[label=\textbf{Q\arabic*.}]
\item Under the read model of Part D, is $S_{\W} = S$? Problem 13 gives only
      $\score(T)\le S_{\W}\le S$. Find a proof, an extra condition that makes
      it true, or a counterexample.
\item The condition $2(1+\lambda)e < \Delta_\lambda$ in Problem 12 assumes the
      worst case for every substitution. Can it be weakened?
\item With $\lambda = 0$ the margin can be negative for our sequences. Is there
      a simple property of the reference that guarantees a positive margin?
\item Extend Problem 12 to reads with two indels.
\end{enumerate}

\section*{Sign-up}

\begin{center}
\renewcommand{\arraystretch}{1.7}
\begin{tabular}{|c|c|c|c|p{0.28\textwidth}|c|c|}
\hline
Part & Problems & Level & Time & Names & Written & Checked by \\ \hline
A & 1--4 & \hard & 2--3 h & & & \\ \hline
B & 5 & \medium & 30 min & & & \\ \hline
C & 6--10 & \medium & 1.5--2 h & & & \\ \hline
D & 11--15 & \hard & 2--3 h & & & \\ \hline
E & 16 & \easy & 20 min & & & \\ \hline
F & 17--18 & \medium & 45 min & & & \\ \hline
G & open & ? & open & & & \\ \hline
\end{tabular}
\end{center}

\end{document}
